% mazzi: 02
\chapter{Il sistema neuromuscolare}
\label{cap:neuromuscolare}

\textbf{Sistema neuromuscolare}: aspetto principale della prestazione, e quindi il presupposto
per capire \textbf{come} e \textbf{che cosa} valutare.

% ---------------------------------------------------------------------------
\section{I fattori limitanti la prestazione sportiva}

La prestazione sportiva si basa sull'\textbf{individualità del singolo atleta}, condizionata da
tre fattori:
\begin{itemize}
  \item \textbf{genotipo}: corredo genetico di partenza, in cui entrano il grado di
    allenabilità e l'età del soggetto, intesa come \textit{età biologica};
  \item \textbf{allenamento}: tipo di lavoro svolto;
  \item \textbf{stato di salute}, a sua volta condizionato da:
    \begin{itemize}
      \item \textbf{infortuni}: nell'arco della carriera il giocatore ha un alto rischio di
        subirne, rischio maggiore per gli eventi successivi al primo;
      \item \textbf{fatica}: incide in maniera considerevole sul sistema neuromuscolare, ed è
        in grado di alterare la prestazione;
      \item \textbf{dieta}.
    \end{itemize}
\end{itemize}

I tre confluiscono sulla prestazione attraverso le leggi della \textbf{fisiologia}, della
\textbf{psicologia} e della \textbf{biomeccanica}: genotipo + allenamento + salute
$\rightarrow$ atleta $\rightarrow$ prestazione sportiva.

% ---------------------------------------------------------------------------
\section{Il sistema nervoso nel controllo del movimento}

Il sistema nervoso interviene nel movimento con quattro sottosistemi:
\begin{itemize}
  \item \textbf{sistema piramidale}: dall'\textbf{area motoria 4} arriva ai muscoli e invia lo
    stimolo per effettuare il movimento; via del movimento volontario;
  \item \textbf{sistema extrapiramidale}: via del movimento automatico e del controllo
    posturale;
  \item \textbf{sistema propriocettivo}: informa le strutture nervose superiori di tutte le
    variazioni che avvengono nel corpo;
  \item \textbf{sistema esterocettivo}: fornisce informazioni sull'ambiente in cui ci si muove.
\end{itemize}

\rubrica{Influenza sul movimento e sulla postura}
\textbf{Movimento} e \textbf{stabilizzazione posturale} dipendono da:
\begin{itemize}
  \item \textbf{performance muscolare}: velocità, forza, potenza, resistenza, accelerazioni e
    decelerazioni;
  \item \textbf{equilibrio posturale};
  \item \textbf{integrità legamentosa e articolare}.
\end{itemize}

Il circuito di controllo procede così:
\begin{enumerate}
  \item \textbf{recettori cutanei}, \textbf{muscolari} e \textbf{articolari} $\rightarrow$
    midollo spinale afferente;
  \item i \textbf{sistemi motori spinali} generano le risposte riflesse $\rightarrow$ midollo
    spinale efferente $\rightarrow$ attività muscolare riflessa;
  \item le informazioni salgono come \textbf{vista e sensibilità esterocettiva},
    \textbf{sensibilità propriocettiva} e \textbf{informazioni labirintiche};
  \item da qui si dividono in due vie:
    \begin{enumerate}
      \item \textbf{sistemi motori corticali} (coscienti) $\rightarrow$ propriocettività
        cosciente, programmazione motoria, tempo di reazione, sincronizzazione neuromuscolare
        $\rightarrow$ controllo della contrazione muscolare volontaria;
      \item \textbf{sistemi motori sopraspinali} (automatici) $\rightarrow$ percezione
        cinestesica cosciente $\rightarrow$ controllo dell'attività muscolare automatica.
    \end{enumerate}
\end{enumerate}

I parametri cinematici che si misurano --- velocità, spostamento, accelerazione, forza, potenza
muscolare, capacità di resistenza --- \textbf{sottendono quindi un'attività del sistema nervoso
a più livelli}, compresa quella dei meccanocettori e dei propriocettori sparsi in muscoli,
tendini, articolazioni e legamenti, che interagiscono con le afferenze visive e vestibolari.
Dietro ai parametri meccanici e biomeccanici misurati ci sono veri e propri \textbf{schemi
neurali}: il punto è decisivo nello studio dell'infortunio, perché qualsiasi infortunio tende ad
alterare il sistema neurale \textbf{in maniera più stabile} di quanto non faccia con le strutture
biomeccaniche.

% ---------------------------------------------------------------------------
\section{La contrazione muscolare}

Le caratteristiche della contrazione sono \textbf{biofisiche}: \textit{biofisico} è il terreno
comune in cui interagiscono le leggi della \textbf{fisica}, che riguardano l'ambiente e il corpo
umano, e quelle della \textbf{biologia}.

Lo stimolo parte dall'\textbf{area motoria 4}, arriva al \textbf{midollo spinale} e da qui,
attraverso l'innervazione dell'\textbf{$\alpha$-motoneurone}, innesca la sequenza che trasforma
l'\textbf{impulso nervoso in energia meccanica}, trasferendo la forza della contrazione al
tendine e quindi all'osso.

\rubrica{Dall'impulso all'accorciamento}
\begin{enumerate}
  \item l'impulso arriva nella \textbf{zona di innervazione} e si propaga come \textbf{potenziale
    d'azione};
  \item attraverso i \textbf{tubuli T} $\rightarrow$ rilascio di $Ca^{2+}$ dal \textbf{reticolo
    sarcoplasmatico};
  \item il $Ca^{2+}$ libera il \textbf{sito attivo dell'actina};
  \item si forma il \textbf{ponte actomiosinico} (\textit{cross-bridge});
  \item scissione di \textbf{ATP} in \textbf{ADP} $\rightarrow$ scorrimento dell'actina sulla
    miosina;
  \item \textbf{accorciamento del sarcomero}; l'accorciamento di tutti i sarcomeri in serie e in
    parallelo produce l'\textbf{accorciamento muscolare}.
\end{enumerate}

In parallelo agiscono i \textbf{sistemi di controllo} --- fusi neuromuscolari e riflesso spinale
in genere --- che inviano informazioni al midollo spinale per intervenire in modo
\textbf{eccitatorio} o \textbf{inibitorio} secondo il recettore coinvolto, e sono quindi in grado
di modulare il movimento.

% ---------------------------------------------------------------------------
\section{L'unità motoria}

\textbf{Unità motoria} (UM): un motoneurone e tutte le fibre muscolari da esso innervate; è la
parte essenziale della contrazione muscolare.

\begin{itemize}
  \item stimolando il motoneurone, tutte le sue fibre si contraggono \textbf{come una sola
    unità};
  \item il numero di fibre per UM varia da poche unità a parecchie migliaia:
    \begin{itemize}
      \item \textbf{UM piccole} $\rightarrow$ movimenti fini e delicati (muscoli extraoculari);
      \item \textbf{UM ampie} $\rightarrow$ muscoli grandi, come i posturali.
    \end{itemize}
\end{itemize}

\subsection{Tipi di unità motorie}

Due tipi fondamentali:
\begin{itemize}
  \item \textbf{soglia di attivazione bassa}: UM piccole, lente e resistenti alla fatica; nei
    \textbf{muscoli posturali}, che restano contratti a lungo;
  \item \textbf{soglia di attivazione alta}: UM grandi, veloci e facilmente affaticabili; nei
    \textbf{muscoli dinamici}, che permettono la prestazione ad alta intensità.
\end{itemize}

\rubrica{Il ruolo dell'\texorpdfstring{$\alpha$}{alfa}-motoneurone}
Le UM si distinguono \textbf{in funzione delle caratteristiche dell'$\alpha$-motoneurone che le
innerva}: tutte e tre le tipologie di fibre hanno la stessa matrice proteica, ed è il
condizionamento imposto dal singolo motoneurone a differenziarle in rosse, intermedie e veloci.

\rubrica{Le tre tipologie}
\begin{enumerate}
  \item \textbf{S} (\textit{slow}) $\rightarrow$ fibre \textbf{SO} (\textit{slow oxidative}),
    dette \textbf{rosse}:
    \begin{itemize}
      \item motoneurone con frequenza di stimolo molto bassa;
      \item metabolismo aerobico $\rightarrow$ difficilmente affaticabili;
      \item scarsa tensione muscolare;
      \item ricche di enzimi, mitocondri e mioglobina $\rightarrow$ colore rosso; abbondanti nei
        muscoli posturali, sempre attivi.
    \end{itemize}
  \item \textbf{FR} (\textit{fast fatigue-resistant}) $\rightarrow$ fibre \textbf{FOG}
    (\textit{fast oxidative-glycolytic}), dette \textbf{intermedie}:
    \begin{itemize}
      \item motoneurone con frequenza d'impulso superiore;
      \item tensione maggiore delle fibre lente;
      \item \textbf{doppio metabolismo}, aerobico e glicolitico.
    \end{itemize}
  \item \textbf{FF} (\textit{fast fatigable}) $\rightarrow$ fibre \textbf{FG} (\textit{fast
    glycolytic}), dette \textbf{veloci}:
    \begin{itemize}
      \item frequenza di scarica elevata degli $\alpha$-motoneuroni;
      \item elevati gradienti di forza: fibre potentissime;
      \item metabolismo glicolitico anaerobico;
      \item facilmente affaticabili.
    \end{itemize}
\end{enumerate}

\subsection{Il reclutamento delle unità motorie}

\textbf{Principio di Henneman} (o principio della dimensione): utilizzando \textbf{carichi
crescenti} e applicando una tensione leggermente superiore al valore d'inerzia, le UM vengono
reclutate in \textbf{ordine crescente di soglia}, dalle più piccole e lente alle più grandi e
veloci.

\begin{enumerate}
  \item carico basso $\rightarrow$ fibre \textbf{lente} (SO);
  \item aumentando il carico $\rightarrow$ fibre \textbf{intermedie} (FOG);
  \item aumentando ancora $\rightarrow$ fibre \textbf{veloci} (FG);
  \item carico intorno all'\textbf{80\% della forza massima} $\rightarrow$ si reclutano tutte le
    fibre del muscolo;
  \item dall'80\% alla \textbf{forza isometrica massima} la tensione viene ulteriormente
    sviluppata \textbf{non} reclutando nuove fibre, ma aumentando la frequenza degli impulsi.
\end{enumerate}

\rubrica{Reclutamento a carico costante}
Se il carico esterno resta costante --- il \textbf{peso del corpo}, quindi un carico
submassimale --- il reclutamento avviene \textbf{in funzione dell'intensità del movimento}:
\begin{itemize}
  \item movimento lento $\rightarrow$ fibre lente;
  \item aumento di intensità e di velocità di contrazione $\rightarrow$ fibre intermedie;
  \item movimenti esplosivi e \textit{sprint} $\rightarrow$ fibre veloci.
\end{itemize}

Ciò \textbf{non} supera il principio di Henneman: a velocità di contrazione altamente elevate i
tempi sono talmente ridotti che le fibre lente non riescono a intervenire nella contrazione.

\rubrica{Intensità dello stimolo e attivazione delle UM}
Al crescere della percentuale di \textit{pool} reclutato aumenta la percentuale di forza tetanica
sviluppata, con la sequenza S $\rightarrow$ FR $\rightarrow$ F$_{(int)}$ $\rightarrow$ FF. I gesti
si collocano lungo questa scala:
\begin{itemize}
  \item \textbf{stazione eretta} e \textbf{cammino}: reclutamento minimo, entro circa il 25\% del
    \textit{pool};
  \item \textbf{corsa}, a 0,6, 1,2, 1,8 e 3\,m/s: reclutamento crescente, fino a poco più del
    50\%;
  \item \textbf{salto}: reclutamento massimo, oltre l'80\% del \textit{pool}, con forza vicina al
    100\% di quella tetanica.
\end{itemize}

% ---------------------------------------------------------------------------
\section{Variabilità delle risposte muscolari}

Un muscolo scheletrico può essere contratto a \textbf{qualsiasi intensità} e con la \textbf{forza
desiderata}. Contrazioni lente o veloci dell'intero muscolo si ottengono agendo su due
meccanismi:
\begin{enumerate}
  \item \textbf{frequenza dello stimolo} inviata alle fibre muscolari;
  \item \textbf{reclutamento}, cioè numero e dimensioni delle UM coinvolte.
\end{enumerate}

Esempio: sollevare una penna e sollevare un tavolo --- stesso gesto, tensioni muscolari
profondamente diverse.

% ---------------------------------------------------------------------------
\section{Tensione interna e tensione esterna}

\begin{itemize}
  \item \textbf{tensione interna}: in un muscolo scheletrico in contrazione, è la tensione
    generata dalle \textbf{miofibrille} all'interno delle fibre; da sola non produce movimento;
  \item \textbf{elementi elastici in serie} (EES): fibre di \textbf{endomisio},
    \textbf{perimisio}, \textbf{epimisio} e \textbf{tendini}, cui la tensione interna viene
    trasferita;
  \item \textbf{tensione esterna}: la tensione degli EES.
\end{itemize}

\rubrica{Comportamento degli EES}
Gli EES si comportano come \textbf{elastici}: inizialmente si allungano facilmente, ma una volta
allungati diventano rigidi e più adatti a trasferire la tensione esterna alla resistenza
applicata. È l'\textbf{irrigidimento} della struttura elastica allungata a permettere il
trasferimento della tensione muscolare all'osso.

\rubrica{Condizione per il movimento}
Il movimento si produce solo se tensione muscolare prodotta $\neq$ resistenza esterna; se le due
sono uguali, non si ha movimento.

% ---------------------------------------------------------------------------
\section{I tipi di contrazione muscolare}

Due regimi:
\begin{itemize}
  \item \textbf{isometrico} (statico): tensione muscolare $\leq$ carico esterno;
  \item \textbf{non-isometrico} (dinamico): si articola in \textbf{concentrico} ed
    \textbf{eccentrico}.
\end{itemize}

In base alla \textbf{strategia} e alla \textbf{velocità di contrazione}, i movimenti dinamici si
distinguono ulteriormente in:
\begin{itemize}
  \item \textbf{isotonico}: muove un carico in maniera costante; la velocità massima si raggiunge
    \textbf{alla fine} della contrazione, non all'inizio;
  \item \textbf{isocinetico}: alcune macchine consentono di lavorare a \textbf{velocità
    costante}, che però non è un movimento naturale.
\end{itemize}

\begin{center}
\begin{forest} classalbero
[Contrazione muscolare
  [Isometrico\\(statico)]
  [Non-isometrico\\(dinamico)
    [Concentrico [Isotonico][Isocinetico]]
    [Eccentrico [Isotonico][Isocinetico]]
  ]
]
\end{forest}
\end{center}

\subsection{Contrazione isometrica}

Il muscolo sviluppa una \textbf{tensione uguale alla resistenza esterna}, trasmessa al tendine
\textbf{senza produrre movimento esterno}.

\begin{itemize}
  \item il \textbf{ventre muscolare si accorcia} comunque, stirando gli elementi elastici in
    serie rappresentati dal tendine;
  \item tensione muscolare = tensione esterna $\rightarrow$ nessun movimento: i capi articolari
    restano nella posizione di origine.
\end{itemize}

\subsection{Contrazione concentrica}

Il muscolo \textbf{si accorcia} sviluppando una tensione \textbf{maggiore della resistenza
esterna} e produce un movimento che avvicina tra loro i capi articolari; la velocità è sempre
positiva.

\begin{itemize}
  \item il \textbf{picco di velocità} viene \textbf{subito dopo il picco di forza}: quando
    compare, si è già nella fase di distensione;
  \item la \textbf{velocità massima} si raggiunge alla fine, quando la forza torna al livello del
    peso corporeo (\textit{body weight});
  \item segue una piccola \textbf{fase di volo}, dopo la quale la velocità decade secondo
    l'accelerazione di gravità.
\end{itemize}

\subsection{Contrazione eccentrica}

Il muscolo si contrae producendo una tensione \textbf{inferiore alla resistenza esterna}: i capi
articolari \textbf{si allontanano} tra loro.

\begin{itemize}
  \item i livelli di forza prodotti in regime eccentrico sono \textbf{superiori alla forza
    massima} espressa in regime isometrico;
  \item si tratta quindi di \textbf{carichi eccessivi}, oltre la forza isometrica massima:
    esercitazione riservata ad atleti molto esperti, che abbiano già utilizzato il carico e
    possiedano una buona tecnica di sollevamento;
  \item impiego soprattutto nell'\textbf{atletica leggera}, dove picchi elevati di forza e
    potenza sono legati alla prestazione; nei \textbf{giochi di squadra} molto meno consigliabile.
\end{itemize}

% ---------------------------------------------------------------------------
\section{Ciclo stiramento-accorciamento e movimento pliometrico}

\textbf{Ciclo stiramento-accorciamento} (\textit{stretch-shortening cycle}, SSC): sequenza in cui
a una fase di stiramento (\textit{stretch}) segue immediatamente una fase di accorciamento
(\textit{shortening}); corrisponde alla \textbf{contrazione eccentrico-concentrica}.

\begin{itemize}
  \item è una \textbf{strategia evolutiva}, il movimento naturale per eccellenza;
  \item nella \textbf{fase eccentrica} muscolo ed elementi elastici in serie accumulano energia
    elastica $\rightarrow$ restituita nella fase concentrica successiva;
  \item prestazione \textbf{molto più efficiente} di qualsiasi altro tipo di contrazione dinamica.
\end{itemize}

\rubrica{Il tempo di accoppiamento}
\textbf{Tempo di accoppiamento}: tempo che indica il passaggio da una fase all'altra del ciclo.
\begin{itemize}
  \item deve essere il \textbf{più breve possibile};
  \item altrimenti l'energia elastica si disperde in calore e il \textbf{riflesso da stiramento
    non avviene} nella fase concentrica successiva.
\end{itemize}

\rubrica{Movimento pliometrico}
\textbf{Movimento pliometrico}: ciclo di contrazione eccentrico-concentrico caratterizzato
dall'\textbf{impatto con il suolo}.
\begin{itemize}
  \item per essere definito \textbf{regime pliometrico} il \textbf{tempo di contatto al suolo}
    deve realizzarsi in tempi brevissimi;
  \item differenza rispetto al \textit{counter movement jump}, dove il tempo di contatto non è
    rilevante: qui è \textbf{nel tempo di contatto} che si esprime tutta l'azione del sistema
    neuromuscolare, e quindi tutto lo sviluppo di potenza.
\end{itemize}

\rubrica{Modello meccanico}
Ruota e corpo rigido che ruota attorno a uno spigolo; \textbf{modello massa-molla}.

% ---------------------------------------------------------------------------
\section{La stiffness muscolare}

\textbf{Stiffness}: in campo fisiologico indica forza, resistenza, densità e rigidità dei
\textbf{tendini} e del \textbf{tessuto connettivo} del muscolo.

\begin{itemize}
  \item maggiore la \textit{stiffness} di questi tessuti $\rightarrow$ maggiore l'energia
    immagazzinabile durante un movimento eccentrico, poi restituita e liberata nella fase
    concentrica;
  \item è la \textbf{capacità reattiva elastica} che un muscolo produce per eseguire contrazioni
    pliometriche subito dopo il \textbf{prestiramento muscolare};
  \item permette di \textbf{restituire velocemente} l'energia elastica accumulata nella
    deformazione avvenuta durante il breve contatto con il terreno.
\end{itemize}

\rubrica{Stiffness e tempo di contatto}
Più breve il tempo di contatto con il terreno, maggiore il grado di \textit{stiffness} del
soggetto e quindi la sua capacità di esercitare forza nel più breve tempo possibile: un vantaggio
determinante per la prestazione. Interviene in tutte le prestazioni di \textbf{corsa} e di
\textbf{salto} --- salto in lungo, \textit{sprint} --- e soprattutto nella \textbf{fase lanciata},
dove il sistema deve sviluppare alta tensione di forza in tempi di contatto sempre più brevi.

\rubrica{Il ruolo del terreno}
La \textit{stiffness} riguarda \textbf{anche il terreno}: allenarla su terreno soffice non ha
senso. Caso classico l'\textbf{allenamento sulla sabbia}, dove la \textit{stiffness} è assente
$\rightarrow$ degrado non solo della prestazione, ma delle stesse caratteristiche neuromuscolari
del calciatore, che non è più in grado di sviluppare tensioni muscolari nel più breve tempo
possibile.

% ---------------------------------------------------------------------------
\section{L'integrazione sensomotoria}

\rubrica{Recettori muscolari}
\begin{itemize}
  \item \textbf{fusi neuromuscolari} (\textit{muscle spindles}): sensibili alla lunghezza; una
    volta stimolati sono \textbf{eccitatori} nei confronti dell'$\alpha$-motoneurone;
  \item \textbf{organi tendinei del Golgi} (GTO, \textit{Golgi tendon organs}):
    \textbf{inibitori} dell'$\alpha$-motoneurone;
  \item \textbf{corpuscoli del Pacini} (\textit{paciniform corpuscles}): sensibili alle
    variazioni anche di velocità;
  \item \textbf{terminazioni nervose libere} (\textit{free nerve endings}): legate ai livelli
    biochimici locali del muscolo.
\end{itemize}

Nel ciclo stiramento-accorciamento l'aumento di prestazione deriva dall'azione congiunta del
\textbf{riuso dell'energia elastica} e dell'azione \textbf{eccitatoria dei fusi}, mentre i
\textbf{GTO} esercitano il controllo inibitorio.

\rubrica{Il riflesso spinale e il \texorpdfstring{$\gamma$}{gamma}-loop}
\begin{enumerate}
  \item il muscolo \textbf{si allunga} nella fase eccentrica $\rightarrow$ i fusi neuromuscolari
    inviano impulsi che innescano un \textbf{riflesso spinale} (\textit{loop});
  \item il riflesso va a innervare ulteriormente l'\textbf{$\alpha$-motoneurone};
  \item a questa attività si somma sempre l'\textbf{attività discendente} del sistema piramidale,
    cioè volontario;
  \item la prestazione ne risulta aumentata attraverso l'\textbf{azione $\gamma$}, o
    \textbf{$\gamma$-loop}.
\end{enumerate}

L'informazione ha un \textbf{effetto immediato} a livello spinale, ma prosegue nelle vie nervose
superiori --- midollo spinale $\rightarrow$ cervelletto $\rightarrow$ area somatosensoriale
$\rightarrow$ area motoria 4 --- per essere integrata nell'attività volontaria. Vi confluiscono
\textbf{recettori muscolari}, \textbf{recettori cinestesici} e i recettori della sensibilità
cutanea: \textbf{corpuscoli di Meissner} (tatto) e \textbf{terminazioni nervose libere} (dolore e
temperatura).

\rubrica{Fatica nel ciclo stiramento-accorciamento}
La \textit{stiffness} si esprime dentro un sistema complesso, in cui all'attività discendente
dell'area motoria si affiancano i \textbf{sistemi riflessi}. Con la \textbf{fatica} l'azione
riflessa si riduce, e con essa la qualità della \textit{stiffness}:
\begin{enumerate}
  \item deterioramento della funzione muscolare;
  \item ridotta tolleranza all'impatto;
  \item perdita del potenziale di energia elastica;
  \item aumento del lavoro durante la fase di spinta (\textit{push-off}).
\end{enumerate}

\textbf{Vie discendenti} e \textbf{midollo spinale}, tramite le afferenze \textbf{Ia},
\textbf{Ib} e di \textbf{gruppo III e IV}, agiscono sul riflesso e sulla \textit{stiffness} di
anca, ginocchio e caviglia $\rightarrow$ performance; il \textbf{danno muscolare}
(\textit{muscle damage}) interviene su questo circuito.
